% eksperymentalne tłumaczenie na włoski

%

\section{Czoworokąty dwuśrodkowe} % lang-pl

\section{Quadrilateri bicentrici} % lang-it

\begin{definition}
Czworokąt, który jest jednocześnie wpisany w pewien okrąg i opisany na innym okręgu, nazywamy dwuśrodkowym. % % lang-pl
\index{czworokąt!dwuśrodkowy}% % lang-pl
Un quadrilatero che è contemporaneamente inscritto in un cerchio e circoscritto a un altro cerchio si chiama bicentrico. % % lang-it
\index{quadrilatero!bicentrico}% % lang-it
\end{definition}

Przykładami takich czworokątów są kwadraty, prostokątne latawce i niektóre równoramienne trapezy. % % lang-pl
\index{kwadrat}% % lang-pl
\index{latawiec}% % lang-pl
\index{trapez!równoramienny}% % lang-pl
Ich pełną charakteryzację można uzyskać przez połączenie warunków \ref{prp_incircle} oraz \ref{prp_excircle}. % % lang-pl
Ale mamy też inne opisy. % % lang-pl
Esempi di tali quadrilateri sono quadrati, aquiloni rettangolari e alcuni trapezi isosceli. % % lang-it
\index{quadrato}% % lang-it
\index{aquilone}% % lang-it
\index{trapezio!isoscele}% % lang-it
La loro completa caratterizzazione può essere ottenuta combinando le condizioni \ref{prp_incircle} e \ref{prp_excircle}. % % lang-it
Ma ci sono anche altre descrizioni. % % lang-it

\begin{proposition}
    Niech $ABCD$ będzie czworokątem opisanym na okręgu $\Gamma$, który dotyka go w punktach $W$ (na odcinku $AB$), $X$ (na $BC$), $Y$ (na $CD$), $Z$ (na $AD$). % lang-pl
    
    Wtedy następujące warunki są równoważne: % lang-pl
    Sia $ABCD$ un quadrilatero circoscritto al cerchio $\Gamma$, che lo tocca nei punti $W$ (su $AB$), $X$ (su $BC$), $Y$ (su $CD$), $Z$ (su $AD$). % lang-it
    
    Allora le seguenti condizioni sono equivalenti: % lang-it
    \begin{itemize}
        \item na czworokącie $ABCD$ można opisać okrąg, % lang-pl
        \item odcinki $WY$ i $XZ$ są prostopadłe, % lang-pl
        \item si può descrivere un cerchio sul quadrilatero $ABCD$, % lang-it
        \item i segmenti $WY$ e $XZ$ sono perpendicolari, % lang-it
        \item $|AW|/|BW| = |DY|/|CY|$,
        \item $|AC|/|BD| = (|AW| + |CY|) / (|BX| + |DZ|)$,
        \item równoległobok Varignona jest prostokątem.% % lang-pl
              \index{równoległobok!Varignona}% % lang-pl
        \item Il parallelogramma di Varignon è un rettangolo.% % lang-it
              \index{parallelogramma!Varignon}% % lang-it
    \end{itemize}
    % TODO: \todofoot{brakujący rysunek}
\end{proposition}

% TODO? https://en.wikipedia.org/wiki/Bicentric_quadrilateral#Construction
% https://en.wikipedia.org/wiki/Bicentric_polygon

Pole powierzchni czworokąta dwuśrodkowego o bokach długości $a, b, c, d$ wynosi $S = \sqrt{abcd}$, jest to prosty wniosek ze wzoru Brahmagupty \ref{brahmagupta_formula}. % % lang-pl
\index{wzór!Brahmagupty}% % lang-pl
Mamy $4r^2 \le S \le 2R^2$, a nawet % % lang-pl
L'area di un quadrilatero bicentrico con lati di lunghezza $a, b, c, d$ è $S = \sqrt{abcd}$, un semplice risultato dalla formula di Brahmagupta \ref{brahmagupta_formula}. % % lang-it
\index{formula!Brahmagupta}% % lang-it
Abbiamo $4r^2 \le S \le 2R^2$, e persino % % lang-it
\begin{equation}
    S \le r^2 \left(1 + \sqrt{\left(\frac{2R}{r}\right)^2 + 1} \right),
\end{equation}
z równością wtedy i tylko wtedy, kiedy czworokąt jest prostokątnym latawcem. % % lang-pl
Inne nierówności, których źródła nie podamy, to % % lang-pl
con uguaglianza se e solo se il quadrilatero è un aquilone rettangolare. % % lang-it
Altre disuguaglianze, le cui fonti non citiamo, sono % % lang-it
\begin{equation}
    2 \sqrt {S} \le p \le r + \sqrt{r^2 + 4R^2},
\end{equation}
gdzie $p$ to połowa obwodu albo % % lang-pl
dove $p$ è metà del perimetro oppure % % lang-it
\begin{equation}
    S \le \frac 1 6 \left(ab + ac + ad + bc + bd + cd\right).
\end{equation}
Nierówność % % lang-pl
La disuguaglianza % % lang-it
\begin{equation}
    R \ge \sqrt 2 r,
\end{equation}
z równością tylko dla kwadratu, jest nietrywialna. Dowiedzie jej Fejes Tóth w 1948 roku. % % lang-pl
con uguaglianza solo per il quadrato, è non banale. Dimostrata da Fejes Tóth nel 1948. % % lang-it
\index[persons]{Tóth, Fejes}%

\begin{theorem}[Fussa, 1792] % lang-pl
Niech $x$ oznacza odległość między środkami okręgu wpisanego i opisanego na czworokącie dwuśrodkowym. % % lang-pl
Wtedy % % lang-pl
\begin{theorem}[Fuss, 1792] % lang-it
Sia $x$ la distanza tra i centri del cerchio inscritto e del cerchio circoscritto del quadrilatero bicentrico. % % lang-it
Allora % % lang-it
\begin{equation}
    \frac{1}{(R-x)^2} + \frac{1}{(R+x)^2} = \frac{1}{r^2}.
\end{equation}
\index{twierdzenie!Fussa}% % lang-pl
\index{teorema!Fuss}% % lang-it
\end{theorem}

% TODO: to jest odpowiednik wzoru Eulera dla trójkątów

Wyznaczając $x$ z twierdzenia Fussa i rozwiązując nierówność $x^2 \ge 0$ dochodzimy znowu do nierówności Tótha. % % lang-pl
Aż dziwne, że nikt tego wcześniej nie zrobił. % % lang-pl
Nicolaus Fuss był szwajcarskim matematykiem, który spędził większość swego życia w Rosji. % % lang-pl
Determinando $x$ dal teorema di Fuss e risolvendo la disuguaglianza $x^2 \ge 0$ si arriva nuovamente alla disuguaglianza di Tóth. % % lang-it
È sorprendente che nessuno l'abbia fatto prima. % % lang-it
Nicolaus Fuss era un matematico svizzero che trascorse gran parte della sua vita in Russia. % % lang-it
\index[persons]{Fuss, Nicolaus}%

\begin{proposition}
Punkt przecięcia przekątnych, środek okręgu wpisanego i środek okręgu opisanego na czworokącie dwuśrodkowym są współliniowe. % % lang-pl
\index{współliniowy}% % lang-pl
Il punto di intersezione delle diagonali, il centro del cerchio inscritto e il centro del cerchio circoscritto del quadrilatero bicentrico sono collineari. % % lang-it
\index{collineare}% % lang-it
\end{proposition}

\begin{proposition}
Niech $ABCD$ będzie czworokątem dwuśrodkowym, zaś $O$ środkiem okręgu opisanego. % % lang-pl
Wtedy środki okręgów wpisanych w trójkąty $\triangle OAB$, $\triangle OBC$, $\triangle OCD$, $\triangle ODA$ leżą na jednym okręgu. % % lang-pl
Sia $ABCD$ un quadrilatero bicentrico e $O$ il centro del cerchio circoscritto. % % lang-it
Allora i centri dei cerchi inscritti nei triangoli $\triangle OAB$, $\triangle OBC$, $\triangle OCD$, $\triangle ODA$ giacciono su un unico cerchio. % % lang-it
\end{proposition}

Wreszcie Klamkin pokazał w 1967 roku, że % % lang-pl
Infine Klamkin dimostrò nel 1967 che % % lang-it
\begin{proposition}
Niech $p, q$ będą długościami przekątnych czworokąta dwuśrodkowego o bokach długości $a$, $b$, $c$, $d$. % % lang-pl
Wtedy % % lang-pl
Siano $p, q$ le lunghezze delle diagonali di un quadrilatero bicentrico con lati di lunghezza $a$, $b$, $c$, $d$. % % lang-it
Allora % % lang-it
\begin{equation}
    8 pq \le (a + b + c + d)^2
\end{equation}
\end{proposition}

Bogdańska, Neugebauer \cite[s. 267]{neugebauer_2018} na ostatniej stronie podają niespodziewanie informacją, że twierdzenie Ponceleta było motywem przewodnim całego skryptu. % % lang-pl
% {\color{red}\textbf{(TODO: T2.19)}\color{black}}
Zachęcają do uogólnienia czwartego dowodu dla poniższej wersji: % % lang-pl
Bogdańska, Neugebauer \cite[p. 267]{neugebauer_2018} all'ultima pagina riportano sorprendentemente che il teorema di Poncelet era il filo conduttore dell'intero manoscritto. % % lang-it
Incoraggiano a generalizzare la quarta dimostrazione per la seguente versione: % % lang-it

\begin{theorem}[Ponceleta, małe] % lang-pl
Niech trójkąt $A_0 A_1 A_2$ będzie wpisany w~stożkową $C$ oraz opisany na stożkowej $D$. % % lang-pl
Wtedy każdy punkt $B_0$ stożkowej $C$ jest wierzchołkiem dokładnie jednego trójkąta $B_0 B_1 B_2$ wpisanego w~stożkową $C$ oraz opisanego na stożkowej $D$. % % lang-pl
\index{twierdzenie!Poneceleta, małe i duże}% % lang-pl
\begin{theorem}[di Poncelet, piccolo] % lang-it
Sia il triangolo $A_0 A_1 A_2$ inscritto nella conica $C$ e circoscritto alla conica $D$. % % lang-it
Allora ogni punto $B_0$ della conica $C$ è il vertice di un unico triangolo $B_0 B_1 B_2$ inscritto in $C$ e circoscritto a $D$. % % lang-it
\index{teorema!Poncelet, piccolo e grande}% % lang-it
\end{theorem}

Jest też wielkie twierdzenie Ponceleta, udowodnione przez, jak niezbyt trudno się domyślić, Victora Ponceleta \cite[s. 311-317]{poncelet_1865} (wg Bogdańskiej, Neugebauera w 1813 roku, wg angielskiej Wikipedii w~1822 roku). % % lang-pl
Esiste anche il grande teorema di Poncelet, dimostrato, come non è difficile intuire, da Victor Poncelet \cite[s. 311-317]{poncelet_1865} (secondo Bogdańska, Neugebauer nel 1813, secondo Wikipedia inglese nel 1822). % % lang-it
\index[persons]{Poncelet, Victor}%

\begin{theorem}[Ponceleta, wielkie] % lang-pl
\begin{theorem}[Poncelet, grande] % lang-it
\label{big_poncelet}%
Niech $C$ i $D$ będą dwiema stożkowymi, zaś $A_0, A_1, \ldots, A_{n-1}$ takimi punktami na stożkowej $C$, że proste $A_0A_1$, $A_1A_2$, \ldots, $A_{n-1}A_0$ są styczne do stożkowej $D$. % % lang-pl
Wtedy dla każdego punktu $B_0$ na stożkowej $C$ istnieją różne punkty $B_1, \ldots, B_{n-1}$, też na stożkowej $C$, że proste $B_0B_1$, $B_1B_2$, \ldots, $B_{n-1}B_0$ są styczne do stożkowej $D$. % % lang-pl
Siano $C$ e $D$ due coniche, e $A_0, A_1, \ldots, A_{n-1}$ punti su $C$ tali che le rette $A_0A_1$, $A_1A_2$, \ldots, $A_{n-1}A_0$ siano tangenti a $D$. % % lang-it
Allora per ogni punto $B_0$ su $C$ esistono punti distinti $B_1, \ldots, B_{n-1}$, anch'essi su $C$, tali che le rette $B_0B_1$, $B_1B_2$, \ldots, $B_{n-1}B_0$ siano tangenti a $D$. % % lang-it
\end{theorem}

Dowód można znaleźć na przykład u Akopiana, Zasławskiego \cite[s. 93, 61, 67, 115, 124]{akopyan_2007}. % % lang-pl
La dimostrazione si può trovare ad esempio in Akopyan, Zaslavsky \cite[p. 93, 61, 67, 115, 124]{akopyan_2007}. % % lang-it

% TODO: https://www.deltami.edu.pl/2012/12/wedrowki-po-okregu/
%